% Paper 3: RML - Introduction Section (LaTeX-ready)
% Meta-Policy Learning Over QA Manifolds
% Generated: 2025-12-29

\section{Introduction}

Learning and control on discrete state spaces governed by algebraic invariants presents a fundamental challenge: how can agents navigate complex manifolds efficiently when ground-truth queries are expensive? Traditional approaches fall into two paradigms. Model-based reinforcement learning learns forward dynamics models $P(s' \mid s, a)$ and plans over simulated trajectories, requiring extensive interaction to capture stochastic transitions. Symbolic planning operates over hand-coded rules and constraints, achieving determinism at the cost of expert knowledge engineering. Both frameworks optimize for task success—maximizing cumulative reward or satisfying goal constraints—but neither directly addresses a more fundamental question: \textit{can learning reduce the cost of querying ground truth?}

We introduce \textbf{structure-aware control}, a paradigm where learned topological predicates guide action selection with minimal oracle usage. Rather than learning ``what happens next'' (dynamics) or ``how much reward to expect'' (value functions), structure-aware policies learn \textit{which worlds are possible}—structural properties of the state space that determine reachability. By querying learned predictions for action scoring and verifying only top choices with the oracle, such policies trade exhaustive ground-truth simulation for efficient structural inference.

This work builds on two prior results. Paper 1 established that Quantum Arithmetic (QA) transition systems—discrete manifolds defined by modular constraints and generator algebras—exhibit rich algebraic structure: 21-element invariant packets, strongly connected component hierarchies, and deterministic failure taxonomies. Paper 2 demonstrated that topological properties of QA manifolds are \textit{learnable from sparse samples}: a World Model (QAWM) trained on 5,000 labeled transitions (1.4\% of Caps(30,30)'s state space) achieved 0.836 AUROC on return-in-$k$ prediction and generalized to 66\% larger manifolds with 0.816 AUROC. Critically, QAWM learns \textit{structural predicates}—binary properties about reachability—not dynamics models.

Here we demonstrate that \textbf{learned structural predictions enable oracle-efficient control}. We define a standard reachability task on Caps(30,30): starting from random off-diagonal states, reach the diagonal target class $\{(b,b) : 1 \le b \le 30\}$ within $k=20$ steps using generators $\Sigma = \{\sigma, \mu, \lambda_2, \nu\}$. We compare three policies:

\begin{enumerate}
\item \textbf{Random-Legal}: Uniform sampling among legal generators (uninformed baseline).
\item \textbf{Oracle-Greedy}: Uses ground-truth BFS to compute return-in-$k$ for each generator, selecting the action maximizing target reachability (information-optimal upper bound).
\item \textbf{QAWM-Greedy}: Scores generators using QAWM's learned return-in-$k$ predictions from Paper 2, then verifies only the top-scoring choice with the oracle (structure-aware policy).
\end{enumerate}

Our central finding is that \textit{learned queries can dominate ground-truth simulation in oracle-limited regimes}. QAWM-Greedy achieves \textbf{4.20 successes per oracle call} versus Oracle-Greedy's 2.97—a 1.41$\times$ efficiency advantage per success—despite lower absolute success rates (32\% vs 60\%). This result proves that when oracle access is expensive, learned topological knowledge provides better efficiency per success than exhaustive ground-truth queries. The trade-off—sacrificing some task success for dramatic oracle reduction—is favorable precisely when structural properties can be learned offline but online interaction is costly.

An ablation study reveals a deeper principle: \textit{topology dominates constraints for planning on discrete manifolds}. Using only QAWM's return-in-$k$ predictions (ignoring legality predictions) achieves 28\% success compared to 10-20\% for scoring functions that combine legality and reachability. This finding suggests that global, multi-step topological structure provides stronger control signals than local, one-step feasibility constraints—a design principle applicable beyond QA to any hierarchically structured state space.

We emphasize that structure-aware control is \textit{not} reinforcement learning. QAWM-Greedy does not learn value functions, does not predict next states, does not optimize cumulative reward, and does not train on task-specific rewards. Instead, it queries pre-trained structural predicates—learned in Paper 2 on generic return-in-$k$ labels—to score actions, then verifies selections with the oracle. This design enables cross-task transfer (the same QAWM model works for different targets), offline training (no online exploration required), and deterministic verification (oracle calls are for validation, not learning).

Our contributions are:

\begin{itemize}
\item \textbf{Oracle efficiency via learned structure}: QAWM-Greedy achieves 1.41$\times$ more successes per oracle call than Oracle-Greedy, demonstrating that learned topological queries dominate ground-truth simulation in resource-constrained settings.
\item \textbf{Topology-over-constraints principle}: Scoring generators by global reachability (return-in-$k$) outperforms combining legality and reachability, revealing that multi-step topological structure dominates one-step feasibility for control.
\item \textbf{Structure-aware paradigm}: We establish a framework—learn topological predicates offline, query for action scoring, verify with oracle—that differs fundamentally from both model-based RL (dynamics simulation) and symbolic planning (hand-coded rules).
\item \textbf{Cross-task generalization}: QAWM-Greedy uses Paper 2's model without task-specific retraining, demonstrating that learned topology transfers to new control objectives.
\end{itemize}

The remainder of the paper is organized as follows. Section 2 reviews related work in model-based RL, symbolic planning, and topology learning. Section 3 describes our experimental setup: the diagonal reachability task, baseline definitions, and evaluation metrics. Section 4 presents our primary results, including the normalized success metric and scoring function ablations. Section 5 discusses theoretical implications: why learned structure dominates simulation, why topology beats constraints, and how structure-aware control differs from RL. We conclude with applications to combinatorial optimization, neural architecture search, and other domains where discrete dynamics exhibit learnable topological structure.
